% Source: 44_Chapter_33_The_Heart_and_Its_Reasons_Pascal_and_th.txt
\bookchapter[The Heart and Its Reasons]{CHAPTER 33}{THE HEART AND ITS REASONS}{Pascal and the Pietists}

\leadin{W}{HAT DOES ONE DO} with the personal belongings of a deceased loved one? It is an all but
universal problem.

A few days after Blaise Pascal’s death, a servant of the house noticed by chance a curious bulge in
the great scientist’s jacket. Opening the lining, he withdrew a folded parchment written in Pascal’s
hand. Within the parchment was a scribbled sheet of paper, a faithful copy of the words of the
parchment:

\begin{poem}
The year of grace 1654.\\
Monday, November 23rd.,~.~.~. from about half past ten in the evening until about half past twelve\\
God of Abraham, God of Isaac, God of Jacob, not of the philosophers and scholars.\\
Certainty, certainty, feeling, joy, peace.\\
God of Jesus Christ~.~.~.\\
I have separated myself from Him, I have fled from Him, renounced Him, crucified Him.\\
May I never be separated from Him~.~.~.\\
Renunciation, total and sweet.
\end{poem}

The words were the record of Pascal’s mystical illumination, two hours in the presence of God! For
eight years he had hidden them in his coat, sewing and unsewing them as he had need. It was his
personal secret, his memorial of the coming of grace to his soul.

Even in the Age of Reason the thirst of the soul could not be ignored. An important new movement
called Pietism arose as a reminder of that fact.

The Roman Catholic Blaise Pascal and the early Pietists, mostly professing Lutherans, would not seem
to belong in the same stream of Christian history. With the hostility of faith prevalent in the
seventeenth and eighteenth centuries, traditional religious labels meant less and less. Pascal’s
powerful critiques of the Enlightenment showed that Enlightenment thinkers could not consistently
defend their narrow technical understanding of reason without major assumptions or leaps of faith.
Pascal’s important witness did not rest in some pretentious, omniscient, self-sufficient version of
human reason, but in a deeper wisdom of reason’s link to the heart. Pascal and the Pietists held to
the heart’s essential role in knowing God. How, then, in the Age of Reason did they sustain the life
of the soul?

\section{APPEALS TO “CHEAP GRACE”}

The Catholic Church of the seventeenth century did not lack spiritual heroes in Europe or in the new
worlds of Asia, Africa, and America. In France, however, she jealously guarded her traditional
powerful alliances. The higher clergy of the church knew their importance in the political order.
They shared with nobles and kings the revenues of the nation and the splendors of the court.

Frequently the teachers of the wealthy and the confessors of the powerful, the Jesuits, were trained
to function in the urbane world of royalty. Their role as spiritual guides of nobles and kings made
the followers of Loyola experts in psychology. They learned to deaden the thunder of Sinai and to
heal the neuroses of remorse by developing the skills of casuistry, the science of right and wrong
for the Christian conscience.

As confessors Catholic priests were expected to know just what was mortal sin or a venial sin or no
sin at all. They had to apply this knowledge and the proper penance to the special experience of
each penitent. In what cases might the letter of the law be set aside for its spirit? When might a
Christian lie or steal or kill for a higher purpose?

Some experts in this science of casuistry were strict interpreters of the moral law; others were
lenient. The Jesuits were lenient—very lenient. They argued that it might be possible to withhold
the truth by “mental reservation” or even tell a lie if some high purpose called for it. They made
so many allowances for sinful human nature that many earnest people protested what seemed to them
“cheap grace,” forgiveness without contrition.

\begin{sidebar}{Mental Reservation and Probabilism}
Wide mental reservation allows the speaker’s ambiguity to deceive the listeners. With Roman soldiers
pursuing Athanasius, his friend told the soldiers that Athanasius was “not far away” when the bishop
was actually hidden in his boat. The soldiers hurried ahead, thinking that they might still overtake
Athanasius’ boat. Strict mental reservation allows the speaker to mislead the listener in light of a
significant qualification that he held in mind. The speaker might reason, “I must not admit that I
told a lie because I had a very noble reason for telling it.” Jesuits reasoned an answer under such
circumstances may not be counted as a lie or a sin.

Jesuits also proposed the theory of probabilism. This theory justified behavior on the grounds that
such behavior had even a remote probability of being morally acceptable.
\end{sidebar}

The most aggressive opposition to the Jesuits came from a movement called Jansenism. Cornelius
Jansen (1585–1638) was a Dutchman who had adopted St. Augustine’s views of sin and grace at the
University of Louvain. He came to believe that the best way to defend Catholicism against the
Calvinist challenge was to return to the doctrines of the great North African and establish a
rigorous moral code for the Catholic clergy to combat the easygoing ethics of the Jesuits.

Jansen carried on his campaign against the Jesuits as professor of Scripture at Louvain and later as
Bishop of Ypres. At his death in 1638 he left, not quite finished, a major treatise, Augustinus,
which soon after its publication in 1640 became the platform for Jansenism in France.

To correct the errors of Catholicism in his own day, Jansen went back to Augustine and argued that
God, before the foundation of the world, had chosen those men and women who should be saved. The
good works of men could never earn salvation without the help of divine grace because man’s will is
not free and his nature is corrupt beyond self-salvation. Only God’s grace, available through
Christ’s death, can save him.

The Catholic Church had never explicitly repudiated these doctrines of Augustine, but she had let
them slip into the background of her teaching. After the Council of Trent the Jesuit defense of free
will magnified man’s good works to the point of cheapening the redeeming work of Christ—at least so
it seemed to Jansen. The Jesuits, he said, made far too much of human reason and woefully little of
a trustful, unquestioning faith in God.

\section{JANSENISM AT PORT-ROYAL}

Jansen’s close friend, Jean DuVergier, abbot of Saint-Cyran, carried the Jansenist message to France
in 1633 when he was appointed confessor of a Cistercian convent called Port-Royal, sixteen miles
from Paris. The youthful abbess of the nunnery, Jacqueline Arnauld, had led her flock in the ways of
Christ totally, yet humbly. The reputation of the convent’s devotion attracted not only other women
but a group of devout laymen called Solitaries who pursued personal holiness in the low marshy
valley surrounding Port-Royal.

Among the notable friends of the convent was Jacqueline’s brother, Antoine Arnauld II, a member of
the Sorbonne, the theological faculty of the University of Paris. After Jansen’s death in 1638,
Arnauld assumed leadership of the Jansenist cause. In 1643 he leveled his theological guns at the
Jesuits. Without naming them he blasted the idea that frequent confession could compensate for
frequent sinning.

The Jesuits could scarcely sidestep the charge. They turned to the pope and impressed upon him the
dangers of Jansenism. Calvinism, they called it, in Catholic garb. In 1653 the pope condemned five
propositions allegedly taken from Augustinus.

Arnauld, however, continued his attack on the Jesuits. He published two Letters to a Duke and Peer
attacking what he claimed were Jesuit methods in the confessional. The Sorbonne now considered a
move to expel him from the faculty.

Arnauld needed a defender, so he turned to a new friend of Port-Royal, a young scientist and master
of French prose, Blaise Pascal (1623–1662). He could not have made a better choice. From childhood
Pascal had known few things other than problems.

The boy’s mother died when he was only three. His father, Stephen Pascal, undertook the education of
the three children, Gilberte, Blaise, and Jacqueline. Occasionally Stephen took his son with him to
the meetings of the Academy of Science and soon the youth’s scientific curiosity was aroused.

Before he reached the age of twenty-seven Pascal had gained the admiration of mathematicians in
Paris, had invented the calculating machine for his father who was a burdened tax collector at the
time, and had discovered the basic principles of atmospheric and hydraulic pressures. He belonged to
the age of scientific greats.

Blaise’s initial contact with the Jansenists came as the result of an accident to his father. One
icy day in January 1646, Stephen hurried out on foot to prevent a duel about to take place. He fell
on the hard frozen ground and dislocated a hip. The two physicians who treated him were devoted
Jansenists. They succeeded not only in curing their patient but also in winning his son to their
doctrines.

They taught the Pascals that physical sufferings were but illustrations of a basic religious truth:
man of himself is a helpless and miserable creature. Blaise had seldom enjoyed a day without pain.
He knew how helpless physicians could be, so the argument struck him with unusual force. It deepened
his sense of the tragic mystery of life.

Blaise also learned from these Jansenist physicians how profoundly the Bible speaks to the human
condition. He became an avid student of Scripture, pondering its pages as he had atmospheric
pressures. The Bible, he felt, was the way to a transformed heart.

In 1651, however, Pascal’s personal tragedy deepened with the death of his father. The loss brought
Blaise to a crisis in his life. His sister, Jacqueline, renounced the world by entering the
Port-Royal convent, and Blaise was left alone in Paris.

He now gave himself to worldly interests. He took a sumptuously furnished home, staffed it with
servants, and drove about town in a coach drawn by four horses. He pursued the ways of “proper
society” as he had once studied geometry. But after a year of pleasure he found only a “great
disgust with the world,” and he plunged into a day-to-day “quiet desperation.” He felt abandoned by
God.

\section{CERTAINTY, JOY, PEACE}

Alone in the darkness of his soul, Blaise turned again to the Bible, to the seventeenth chapter of
the gospel of John where Jesus prepares for his sacrifice on the cross. All of a sudden the Fire
appeared, as he described in his scribbled memorial: “Certainty, certainty, feeling, joy, peace.”

Pascal’s new faith drew him magnetically into the orbit of Port-Royal. Late in 1654 he joined his
sister, Jacqueline, as a member of the community. That is where Arnauld found him when he needed an
advocate in his case against the Jesuits.

Pascal responded brilliantly. He penned eighteen Provincial Letters exposing the Jesuit theology and
practices by flashes of eloquence and sarcastic wit. As each letter appeared, the public snatched
them up. They were instant best sellers. Port-Royal was no longer an obscure monastery; it was the
center of public interest. The pope condemned the letters, predictably, but all educated France read
them, as succeeding generations did for two centuries.

After writing the last of the letters in March 1657, Pascal looked forward to a projected book on
the evidences for Christianity. He was never able to complete the manuscript. In June 1662 he was
seized with a violent illness, and after lingering two months he died on August 19 at the age of
thirty-nine.

Friends found portions of his writing on faith and reason, and eight years after his death they
published these notes under the title Thoughts (Pensées). In these Pascal is the religious genius
who cuts across doctrine and pierces to the heart of man’s moral problem. He appeals to the
intellect by his passion for truth and arouses the emotions by his merciless descriptions of the
plight of man without God.

In the Age of Reason Pascal insisted that nature is no sure clue to God. “I look on all sides,” says
Pascal,

\begin{blockquote}
and everywhere I see nothing but obscurity~.~.~.~. If I saw no signs of a divinity, I would fix
myself in denial. If I saw everywhere the marks of a Creator, I would repose peacefully in faith.
But seeing too much to deny (Him), and too little to assure me, I am in a pitiful state, and I would
wish a hundred times that if a God sustains nature it would reveal Him without ambiguity.
\end{blockquote}

Man, caught in this ambiguous universe, is himself the greatest mystery. He is part angel and part
beast. He reminds Pascal of the Chimera, which, in Greek mythology, was a she-goat with a lion’s
head and a serpent’s tail: “What a Chimera is man! What a novelty, a monster, a chaos, a
contradiction, a prodigy! The glory and refuse of the universe. Who shall unravel this confusion?”

Even reason is no sure guide. If we trust reason alone, we will doubt everything except pain and
death. But our hearts tell us this cannot be true. That would be the greatest of all blasphemies to
think that life and the universe have no meaning. God and the meaning of life must be felt by the
heart rather than by reason: “The heart has its reasons which reason does not know.”

Pascal saw the human condition so deeply yet so clearly that men and women in our own time, after
three centuries, still gain perspective from him for their own mysterious pilgrimage.

After Pascal’s death, the combined opposition of the Catholic Church and King Louis XIV succeeded in
forcing Jansenism out of France. Port-Royal was destroyed and the movement forced to take refuge in
Holland.

In the seventeenth century, however, Roman Catholic France was not the only area in need of heart.
Any religion that becomes the religion of the majority and slowly turns into a social habit tends to
grow humdrum and flat, regardless of its original glow of enthusiasm. So it proved in many areas of
Lutheran Germany.

\begin{sidebar}{Protestant Scholasticism}
During the one hundred years after Luther and Calvin taught, a substantial change came over the
Protestant movement. The dynamic character of early Protestant calls for believers to entrust
themselves to God’s transforming mercy was eclipsed by doctrinal discussions that seemed abstract
and lifeless. Luther had read the Bible to discern its life-giving witness to Jesus. The Protestant
Scholastics fashioned involved, tedious arguments that depended upon Aristotle’s logic (which Luther
had condemned) and not upon Scripture. When later Protestants did appeal to the Bible, they often
abstracted statements as if they were atomic, logical bits of data and not part of the great
historical flow of God’s story. Scripture seemed to be viewed as a self-sufficient body of material
and not the book that came alive by the witness of God’s Spirit. With God seen in less personal
terms, Scripture seemed different: a source, but not a life-changing witness. Also, theology was
shaped by the schools and less by the church. In time, even sermons seemed like lectures or
theoretical ponderings. What Luther and Calvin said mattered less than what they would have said if
they had addressed the speculative topics of Protestant Scholastics. Sadly, this theological
undertaking did not distinguish them from unbelievers or even Catholics, but from one another.
Confessions and creeds divided the church more and more.
\end{sidebar}

On the heels of the vigorous, creative movement called the Reformation came a cautious period called
Protestant Scholasticism or confessionalism. Out of the depths of his own experience, Luther had
proclaimed a robust doctrine of justifying faith. In the seventeenth century, however, his dedicated
followers, under the spell of the intellect, turned faith into a mental exercise. No longer an act
of surrender to the mercy of God revealed in Christ, faith was now a formal assent to doctrinal
truths set forth by scholars.

The Christian life was less a personal relation to Christ and more a matter of membership in the
state church. Faithful attendance at public worship and reception of the sacraments offered by
orthodox ministers were the essential marks of a good Christian.

Pietism arose as a reaction to this ossification of the Reformation. Just as Jansenism opposed the
cheap grace of the French Jesuits, so the Pietists challenged the nominal faith of German
Lutheranism.

The aims of the Pietists were twofold. First, they stressed the importance of personal faith. They
left behind all dreams of Catholic Christendom and Puritan commonwealths. They believed that
Christianity started with the individual. So, for the first time in Christian history, the idea of
conversions of baptized Christians (as well as pagans) came to prominence. The essence of faith,
said the Pietist, is a personal experience of God’s grace in the believer’s heart.

Second, the Pietists wanted to shift the center of the Christian life from the state churches, in
which a person was born and brought up, to intimate fellowships of those who had a living faith in
God. Revitalized laymen from these centers were expected to spread the Word of God through all
classes of men.

Three men—Philip Jacob Spener, August Hermann Francke, and Count Nikolaus von Zinzendorf—highlight
the history of Pietism.

\section{THE SOUL’S RETURN TO LIFE}

Philip Spener (1635–1705) grew up under strong religious influences, including the writings of the
German mystic Johann Arndt and the English Puritans. Later, at the University of Strasbourg, he met
professors who introduced him to Luther and understood justification by faith not simply as a
doctrine but as a spiritual rebirth.

After three years in a pastorate at Strasbourg, Spener accepted a call to the important Lutheran
city of Frankfort. He was shocked by conditions in the town. He abandoned the prescribed texts and
began to preach from the whole Bible. He called for repentance and discipleship. For several years
nothing earth-shaking happened. Then in 1669 he preached from the Sermon on the Mount and response
was sudden and surprising. People were converted and family life changed.

Spener gathered a little company of dedicated believers in his house twice weekly for reading of
Scripture and religious conversation. These meetings were soon called, in scorn, “gatherings of the
pious,” and Pietism was born.

As interest in devotional literature grew, Johann Arndt’s sermons were published and Spener wrote an
introduction to them. He called it Pious Desires. In it he recommended the establishment of Bible
study groups for spiritual development; a strenuous, rather ascetic, Christian life; greater care
for the Christian character of theological students; and simpler and more spiritual preaching.
Spener’s thought was that a cell of experiential Christians should be gathered in each congregation
to cultivate a stricter and warmer Christian life. He hoped that this leaven would permeate the
whole church. He felt that only those who had been born again by a conscious Christian experience,
conversion, were fitted for this work or should have a place in the ministry.

In 1686 Spener accepted a call to become court preacher at Dresden. But because of his
uncompromising preaching he was often in trouble with the authorities, and in 1692 he welcomed an
invitation from the elector of Brandenburg to move on to Berlin. That same year he persuaded
Frederick, the future king of Prussia, to invite August Hermann Francke (1663–1727) to become a
professor at the new University of Halle. In this Spener showed great wisdom and humility. Francke
soon rose to leadership of the Pietist movement, though Spener continued writing and preaching until
his death in February 1705.

Francke’s arrival in Halle followed several years at the University of Leipzig, where his spiritual
emphasis had created tensions with his faculty colleagues.

\begin{blockquote}
He traced his zeal to a life changing conversion that had come after two years of inner struggle and
doubt. While writing a sermon on John 20:31, Francke had dropped to his knees in great fear. “I
implored,” he said, “the God, whom I did not yet know and whom I did not believe, that, if there
really is a God, he should save me from this wretched condition.” God heard him! Sadness left his
heart and he was “suddenly overwhelmed by a flood of joy.” When he arose, he knew he was saved and
from that time on, it was “easy to live righteously and joyfully in this world.”
\end{blockquote}

At Halle, Francke pioneered an array of spiritual and social ministries. The university became the
hub of a host of Pietist ministries. Francke’s compassion for the neglected led him to begin a
school for the poor. He also established an orphanage and bought a tavern and adjunct land to build
a hospital. His ongoing work included a Latin school for talented boys, a house for widows, a house
for unmarried women, a medical dispensary, a book depot, a printing establishment, and a Bible
house. When King Frederick IV of Denmark wished to establish one of the earliest of Protestant
missions in India in 1705, it was among Francke’s disciples in Halle that he found his first
missionaries.

\section{THE CARPENTER AND THE COUNT}

The later phase of Pietism is dominated by Count von Zinzendorf (1700–1760), an ardent and emotional
man who believed that the mark of true Christianity is a simple, childlike faith in the blood of
Jesus. In vivid, almost erotic images he sang of the soul’s relation to Christ.

Zinzendorf was the grandson of a nobleman who had left Austria for his religious convictions. When
young Nikolaus’s father died and his mother remarried, his training fell into the hands of his
devout grandmother who was an admirer of the Halle Pietists. Nikolaus studied law for three years at
the University of Wittenberg, but the ministry remained his soul’s concern.

An opportunity for service came in an unexpected way with the remains of the old Hussite movement,
the Bohemian Brethren (Unitas Fratrum). The brethren flourished in Bohemia and Moravia at the time
of the Reformation but had been nearly crushed during the Thirty Years’ War, and were subject to
severe persecution. Under a Moravian carpenter, Christian David, the Brethren had experienced the
stirring of revival and were casting about for some place of refuge in Protestant lands. They found
it on Zinzendorf’s estates, and in 1722 David established a community called Herrnhut, “The Lord’s
Watch.” Zinzendorf entered into the community and soon devout people of various backgrounds were
flocking to Herrnhut. The Moravians had in mind a town inhabited only by Christians, separate from
the World, a real “communion of saints.” It was a free and social monasticism, without celibacy. But
like monasticism, they sought to live the Christian life under peculiarly favorable conditions and
apart from grosser temptations.

From 1727 Zinzendorf was the guiding spirit of Herrnhut, and ten years later he received formal
ordination in the reorganized Moravian Church, or United Brethren as the believers preferred to call
it. Zinzendorf’s impulses were always strongly missionary. As a result the Moravians became the
first large-scale Protestant missionary force in history.

On a visit to the Danish capital, Copenhagen, for the coronation of King Christian VI in 1731,
Zinzendorf met a black man from the Danish West India Islands and was impressed with the needs of
the enslaved people. As a result, in 1732 the first of an army of Moravian missionaries, Leonhard
Dober and David Nitschmann, set forth from Herrnhut on their way to St. Thomas. The establishment of
other missions speedily followed: Greenland, Lapland, the American Indians in Georgia, the African
Guinea Coast, the Hottentots of South Africa, South American Guiana, Ceylon, and Algeria.

Zinzendorf’s character, like most, had its light and its shadow. He was inclined to a sentimental
type of religion. Yet some of his emotional hymns, such as “Jesus still lead on, till our rest be
won,” have found a place in the worship of many churches. Few men have shown such personal devotion
to Christ, and he gave the true cornerstone of his character in his declaration to his Herrnhut
congregation: “I have only one passion. It is He, none but He.”

\section{VIEW FROM A REARVIEW MIRROR}

Pietism made an enormous contribution not only to the German people but to Christianity worldwide.
It shifted emphasis in eighteenth-century churches from avid controversy to the care of souls. It
made preaching and pastoral visitation central concerns of the Protestant ministry. It enriched
Christian music enormously. And it underscored the importance of a spiritual laity for a revived
church. Perhaps its greatest legacy was its emphasis on small groups and the devotional reading of
Scripture. Church meetings did not offer enough systematic exposure to Scripture, while small groups
focused attention on supplementing and encouraging personal Scripture reading and corporate worship.
Personal copies of Scripture were more accessible to believers than ever before in history. Personal
engagement with the Bible became a focal point of discipleship. Groups would highlight a word or
thought from Scripture. A well-chosen hymn would reinforce the daily meditation from Scripture.

Supporting all these emphases was the Pietist’s dominant theme: regeneration. And they meant by this
not a theological doctrine but the indispensable experience of the Christian. They believed this
spiritual rebirth was the true fulfillment of the Protestant Reformation. In the heart Christian
doctrine became Christian reality.

The intensely personal way the Pietist described regeneration often made Christianity a drama of the
human soul. The heart of man was the scene of a desperate struggle between the powers of good and
evil.

In this sense Pietism was the fountain of all modern revivals. It set the experience of new life in
Christ at the center of the Christian message and the Christian ministry. For this reason it is
impossible to think of evangelical Christianity today without the imprint of Pietism.

Evangelicals inherited two important traits from Pietism. First, emotion played so large a part in
the Pietist’s religious life that reason was endangered. Since the mind of man could not fathom the
mysteries of human destiny, feelings were left to carry the meaning of faith. Consequently Pietism
had little to say about God’s place in nature or human history. And it presented few challenges to
the spread of secularism. Evangelicals often surrendered to the same weakness.

Second, Pietism assumed the existence of the institutional church. It made no frontal attacks upon
it. But it shifted what was essential to Christianity—the new birth and the spiritual life—from the
traditional state churches to intimate fellowship groups or voluntary associations of believers. In
the United States, where no state church ever prevailed, Evangelicals made an alliance with the
denominational concept of the church and simply multiplied religious movements stressing evangelism
or some aspect of the Christian life.

Clearly, then, Pietism lives on, not only in the Moravian denomination, but also in evangelical
Christianity at large, the spiritual descendants of John Wesley and George Whitefield.

\readinghead{Suggestions for Further Reading}
\begin{readinglist}
\entry Brown, Dale. Understanding Pietism. Revised Edition. Nappanee, IN: Evangel Publishing House,
  1996.
\entry Cailliet, Emile. Pascal: The Emergence of Genius. New York: Harper \& Brothers, 1961.
\entry Groothuis, Douglas. On Pascal. Melbourne: Thomson/Wadsworth, 2003.
\entry *Hammond, James. The Cambridge Companion to Pascal. Cambridge: New York, 2003.
\entry McGiffert, A. C. Protestant Thought Before Kant. New York: Harper \& Brothers, 1961.
\entry McNeill, John T. Modern Christian Movements. New York: Harper \& Row, 1954.
\entry *O’Connell, Marvin. Blaise Pascal: Reasons of the Heart. Grand Rapids: Eerdmans, 1997.
\entry *Shantz, Douglas H. An Introduction to German Pietism: Protestant Renewal at the Dawn of
  Modern Europe. Baltimore, MD: Johns Hopkins, 2013.
\entry Stoeffler, F. Ernest. The Rise of Evangelical Pietism. Leiden: E. J. Brill, 1965.
\end{readinglist}
